\chapter{计算架构与调度}
\label{chap:compute-architecture-and-scheduling}

\setcounter{footnote}{0}

第一章介绍过，CPU 的设计目标是尽量降低指令执行延迟，而 GPU 的设计目标是尽量
提高指令执行吞吐量。第二章和第三章介绍了 CUDA 编程的核心特征：编写和调用内核，
启动并执行由大量线程组成的网格。在接下来的三章中，我们将讨论现代 GPU 的架构，
包括计算架构、内存架构，以及由理解这些架构而产生的性能优化技术。

本章介绍 CUDA C++ 程序员需要理解、并据此分析内核代码性能行为的若干 GPU 计算
架构特征。我们首先给出计算架构的高层次简化视图，并研究灵活的资源分配、线程块
调度和占用率（occupancy）等概念；随后进一步讨论线程调度、延迟容忍、控制流分歧
和同步。最后，我们介绍可用于查询 GPU 可用资源的 API 函数，以及帮助估计 GPU
执行内核时占用率的工具。

接下来的两章将介绍 GPU 内存架构的核心概念和编程注意事项。第五章重点讨论片上
内存架构；第六章将简要介绍片外内存架构，然后详细讨论整个 GPU 架构的各种性能
考量。掌握这些概念的 CUDA C++ 程序员，将能够编写和理解高性能并行内核。

\section{现代 GPU 的架构}
\label{sec:modern-gpu-architecture}

图 \ref{fig:4-1} 展示了 CUDA C++ 程序员所看到的、典型 CUDA-capable GPU 架构的
高层次视图。GPU 由一组高度多线程化的流式多处理器（Streaming Multiprocessor，
SM）组成。每个 SM 都有多个称为流式处理器（streaming processor）的处理单元
\footnote{流式处理器有时也称为 CUDA core。为避免与 CPU core 混淆，本书统一使用
“流式处理器”这一术语。}，
这些处理单元共享控制逻辑和内存资源。图 4.1 中，SM 内的小方格表示流式处理器。
例如，Hopper H100 GPU 有 132 个 SM，每个 SM 有 128 个流式处理器，因此整个 GPU
共有
\[
  132\times128=16{,}896
\]
个流式处理器。从 Hopper 架构开始，GPU 中的 SM 被分组为 GPU 处理集群（GPU
Processing Cluster，GPC）。例如，Hopper H100 的 132 个 SM 分布在 8 个 GPC 中。

SM 的内存资源由不同的片上内存结构组成，在图 \ref{fig:4-1} 中统称为
\textbf{Memory}。第五章将讨论这些片上内存结构。GPU 还配有若干 GB 的片外设备
内存，在图中称为\textbf{全局内存（Global Memory）}。全局内存通过多个内存控制器
或内存通道访问；这些控制器或通道在整体上提供较高的数据访问带宽。

较早的 GPU 使用 GDDR（Graphics Double Data Rate）SDRAM（Synchronous Dynamic
Random Access Memory）。从 NVIDIA Pascal 架构开始，GPU 还可能使用 HBM
（High-Bandwidth Memory）、HBM2、HBM2E 或 HBM3。这些内存由与 GPU 集成在同一
封装中的 DRAM 模块构成。为简洁起见，本书后文将这些内存类型统称为 DRAM。第六章
将讨论访问 GPU DRAM 时涉及的最重要概念。

\begin{figure}[htbp]
  \centering
  \scriptsize
  \begin{tabular}{c}
    \fbox{\begin{minipage}{0.88\textwidth}
      \centering
      \textbf{GPU}\\[0.3em]
      \begin{tabular}{c@{\hspace{0.3em}}c@{\hspace{0.3em}}c@{\hspace{0.3em}}c}
        \fbox{\begin{minipage}{0.18\textwidth}\centering
          \textbf{GPC}\\[0.2em]
          \fbox{\begin{minipage}{0.82\textwidth}\centering
            \textbf{SM}\\[-0.1em]
            \scriptsize SP SP SP\\[-0.1em]
            \scriptsize SP SP SP
          \end{minipage}}\\[-0.1em]
          \scriptsize $\cdots$ 多个 SM
        \end{minipage}}
        &
        \fbox{\begin{minipage}{0.18\textwidth}\centering
          \textbf{GPC}\\[0.2em]
          \fbox{\begin{minipage}{0.82\textwidth}\centering
            \textbf{SM}\\[-0.1em]
            \scriptsize SP SP SP\\[-0.1em]
            \scriptsize SP SP SP
          \end{minipage}}\\[-0.1em]
          \scriptsize $\cdots$ 多个 SM
        \end{minipage}}
        &
        \fbox{\begin{minipage}{0.18\textwidth}\centering
          \textbf{GPC}\\[0.2em]
          \fbox{\begin{minipage}{0.82\textwidth}\centering
            \textbf{SM}\\[-0.1em]
            \scriptsize SP SP SP\\[-0.1em]
            \scriptsize SP SP SP
          \end{minipage}}\\[-0.1em]
          \scriptsize $\cdots$ 多个 SM
        \end{minipage}}
        &
        \fbox{\begin{minipage}{0.18\textwidth}\centering
          \textbf{GPC}\\[0.2em]
          \fbox{\begin{minipage}{0.82\textwidth}\centering
            \textbf{SM}\\[-0.1em]
            \scriptsize SP SP SP\\[-0.1em]
            \scriptsize SP SP SP
          \end{minipage}}\\[-0.1em]
          \scriptsize $\cdots$ 多个 SM
        \end{minipage}}
      \end{tabular}\\[0.5em]
      \fbox{\parbox[c][0.55cm][c]{0.72\textwidth}{\centering
        \textbf{片外全局内存（DRAM）}\quad
        内存控制器 / 内存通道}}
    \end{minipage}}
  \end{tabular}
  \caption{CUDA-capable GPU 的架构（依据原书图 4.1 重绘的概念示意；
  SP 表示流式处理器）}
  \label{fig:4-1}
\end{figure}

\section{线程块调度}
\label{sec:thread-block-scheduling}

调用内核时，CUDA 运行时系统会启动一个执行内核代码的线程网格。虽然 GPU 拥有
大量执行资源，但一个内核所启动的线程网格很容易需要远超 GPU 当前可用的执行
资源。因此，需要在线程块执行资源变得可用时，将线程分配给这些资源的线程块调度
机制。

在 CUDA 中，线程以线程块为单位分配给 SM。也就是说，一个线程块中的所有线程会
同时分配到同一个 SM。图 \ref{fig:4-2} 展示了这种分配方式。同一个 SM 可以同时
分配多个线程块，具体数量取决于 SM 的可用资源。例如，图中每个 SM 分配了三个
线程块。不过，线程块需要预留执行所需的硬件资源，所以一个 SM 能同时分配的线程
块数量是有限的；这一数量取决于多种因素，第 \ref{sec:resource-partitioning}
节将讨论这些因素。

\begin{figure}[htbp]
  \centering
  \scriptsize
  \begin{tabular}{c@{\hspace{0.7em}}c@{\hspace{0.7em}}c}
    \fbox{\begin{minipage}{0.24\textwidth}\centering
      \textbf{网格}\\[0.2em]
      块 0、块 1、块 2、$\ldots$
    \end{minipage}}
    & $\longrightarrow$ &
    \fbox{\begin{minipage}{0.48\textwidth}\centering
      \textbf{GPU}\\[0.2em]
      \begin{tabular}{cc}
        \fbox{\begin{minipage}{0.18\textwidth}\centering
          \textbf{SM 0}\\[-0.1em]块 0\\块 3\\块 6
        \end{minipage}}
        &
        \fbox{\begin{minipage}{0.18\textwidth}\centering
          \textbf{SM 1}\\[-0.1em]块 1\\块 4\\块 7
        \end{minipage}}\\[0.3em]
        \fbox{\begin{minipage}{0.18\textwidth}\centering
          \textbf{SM 2}\\[-0.1em]块 2\\块 5\\块 8
        \end{minipage}}
        &
        \fbox{\begin{minipage}{0.18\textwidth}\centering
          \textbf{待执行列表}\\[-0.1em]块 9、块 10、$\ldots$
        \end{minipage}}
      \end{tabular}
    \end{minipage}}
  \end{tabular}
  \caption{线程块按块分配给流式多处理器（依据原书图 4.2 重绘）}
  \label{fig:4-2}
\end{figure}

由于 SM 数量有限，而且每个 SM 能同时分配的线程块数量也有限，因此一个 CUDA
设备能够同时执行的线程块总数存在上限。但是，网格通常包含远多于该上限的线程块。
为了确保网格中的所有线程块都得到执行，运行时系统会维护待执行线程块列表，并在
已分配线程块执行完毕后，把新的线程块分配给 SM。

以线程块为单位把线程分配给 SM，可以保证同一个线程块中的线程在同一个 SM 上、
并且在相近的时间执行。这一保证使同一线程块中的线程能够以不同线程块之间无法
实现的方式交互。例如，它们可以进行线程块范围的屏障同步，也可以通过位于 SM
中的低延迟 SRAM（Static Random Access Memory）共享内存交换数据。第五章将讨论
共享内存。

从 Hopper 架构开始，CUDA 引入了一个可选的层次结构：线程块集群（thread block
cluster），即一组线程块。类似于同一线程块中的线程会被分配到同一个 SM，同一
线程块集群中的线程块会共同调度到同一个 GPC。集群中的线程可以通过相应 API
进行同步，并且可以访问由集群中所有线程块共享内存组成的分布式共享内存
（distributed shared memory）。第五章将进一步讨论分布式共享内存。

\section{同步与透明可扩展性}
\label{sec:synchronization-transparent-scalability}

CUDA 允许同一线程块中的线程使用线程块范围的屏障同步函数协调活动。最常用的
函数是 \code{__syncthreads()}。注意，该函数名中的两个下划线分别由两个连续的
\code{_} 字符组成；这是 CUDA 中表示内建函数（intrinsic function）的约定，见
“内建函数”侧栏。

当一个线程调用 \code{__syncthreads()} 时，它会停留在调用所在的程序位置，直到
同一线程块中的每个线程都到达该位置。调用同步函数的线程被称为到达了与该程序
位置对应的屏障。屏障确保线程块中的所有线程都完成当前执行阶段，然后任何线程
才能进入下一阶段。只有当线程块中的所有线程都到达屏障后，它们才会被释放，继续
执行屏障之后的语句。

\begin{quote}
\textbf{内建函数}

现代处理器通常提供能够执行关键功能或显著提高性能的特殊指令。这些指令通常以
内建函数（intrinsic，简称 intrinsic）的形式提供给程序员。从程序员的角度看，
它们像库函数；但编译器会特殊处理它们，把每次调用转换为相应的特殊指令。最终
代码中通常没有普通的函数调用，只有与用户代码内联的特殊指令。GNU Compiler
Collection（gcc）、Intel C Compiler 和 Clang/LLVM C Compiler 等主流现代编译器
都支持内建函数。
\end{quote}

在线程块集群层次，CUDA 通过 Cooperative Groups API \cite{cuda-guide-4} 提供集群
范围同步机制。线程块集群中的所有线程可以在某个程序位置同步，以确保同一集群
中的所有线程都到达屏障，并且它们对分布式共享内存的操作都已完成，然后再从屏障
中释放。参与某个屏障的线程集合称为该屏障的\textbf{范围（scope）}。

在网格层次，整个网格中的线程也可以通过 Cooperative Groups API \cite{cuda-guide-4}
进行网格范围的屏障同步。不过，网格范围同步的开销高于线程块范围和集群范围
同步，并且必须满足若干重要限制，以确保网格中的所有线程确实同时在执行并能够
到达屏障。下面主要讨论线程块范围的屏障同步。

\begin{figure}[htbp]
  \centering
  \scriptsize
  \begin{tabular}{c|l}
    \toprule
    线程 & 执行时间线（从左向右）\\
    \midrule
    $T_0$ & \rule{0.16\textwidth}{0.4pt}\hspace{0.04\textwidth}
      \texttt{barrier}\hspace{0.08\textwidth}\rule{0.20\textwidth}{0.4pt}\\
    $T_1$ & \rule{0.22\textwidth}{0.4pt}\hspace{0.04\textwidth}
      \texttt{barrier}\hspace{0.08\textwidth}\rule{0.20\textwidth}{0.4pt}\\
    $T_2$ & \rule{0.30\textwidth}{0.4pt}\hspace{0.04\textwidth}
      \texttt{barrier}\hspace{0.08\textwidth}\rule{0.20\textwidth}{0.4pt}\\
    $T_N$ & \rule{0.38\textwidth}{0.4pt}\hspace{0.04\textwidth}
      \texttt{last arrival}\hspace{0.04\textwidth}\rule{0.20\textwidth}{0.4pt}\\
    \bottomrule
  \end{tabular}\\[-0.1em]
  \scriptsize 左侧曲线表示到达时间；最后一个线程到达后，所有线程才能继续。
  \caption{屏障同步的执行时序示例（依据原书图 4.3 重绘）}
  \label{fig:4-3}
\end{figure}

屏障同步是协调并行活动的简单而常用的方法。现实生活中也经常需要屏障同步。
例如，四位朋友乘车去购物中心。他们可以分别去不同商店购买自己的物品，这比
四人始终聚在一起、依次访问所有商店高效得多。但是，在离开购物中心之前需要
进行屏障同步：必须等四位朋友都回到车上才能出发。提前完成购物的人必须等待
尚未完成的人；如果没有屏障，汽车可能在某人仍在商场时离开。

图 \ref{fig:4-3} 中有 $N$ 个参与同步的线程，时间从左向右推进。有些线程较早
到达屏障，有些线程较晚到达。早到达的线程等待晚到达的线程；当最后一个线程
到达屏障时，所有线程都可以被释放并继续执行。屏障同步保证“没有人掉队”。

在 CUDA 中，如果存在 \code{__syncthreads()} 语句，则线程块中的所有线程都必须
执行它。当 \code{__syncthreads()} 放在 \code{if} 语句中时，线程块中的所有线程
要么都执行包含该调用的路径，要么都不执行。对于 \code{if-else} 语句，如果两个
路径各有一个 \code{__syncthreads()}，线程块中的所有线程也必须全部执行
\code{then} 路径，或全部执行 \code{else} 路径。两个调用是不同的屏障同步点。

图 \ref{fig:4-4} 展示了一个错误用法：从第 04 行开始的 \code{if} 语句中有两个
屏障同步调用。偶数 \code{threadIdx.x} 的线程执行 \code{then} 路径，其余线程
执行 \code{else} 路径。第 06 行和第 10 行的调用定义了两个不同的屏障。由于不能
保证线程块中的所有线程会到达其中任一屏障，这段代码违反了屏障同步函数的使用规则，
执行行为未定义。错误的屏障同步用法可能导致错误结果，也可能让线程永远相互等待，
即死锁。程序员必须避免这种不恰当的用法。

\begin{figure}[htbp]
  \centering
  \begin{lstlisting}[language=C++,firstnumber=1]
__global__ void foo_kernel(...) {
    ...
    ...
    if (threadIdx.x % 2 == 0) {
        ...
        __syncthreads();
    } else {
        ...
        ...
        __syncthreads();
    }
}
  \end{lstlisting}
  \caption{错误使用 \code{__syncthreads()} 的示例；第 06 行和第 10 行
  是不同的屏障（依据原书图 4.4 重绘）}
  \label{fig:4-4}
\end{figure}

屏障同步会对其范围内的线程施加执行约束。这些线程应当在相近的时间完成，以免
等待时间过长。更重要的是，系统必须确保参与屏障的所有线程都能获得最终到达
屏障所需的资源，否则永远无法到达的线程可能导致死锁。CUDA 运行时通过以线程块
为单位分配执行资源来满足这一约束：线程块中的所有线程不仅必须被分配到同一个
SM，还必须同时被分配到该 SM。只有当运行时为线程块中的所有线程取得所需资源后，
线程块才能开始执行。

这说明 CUDA 屏障同步应尽量限制在较小范围内。如果屏障同步只限于线程块范围，
CUDA 运行时就可以按照任意相对顺序执行不同线程块。这种灵活性使 CUDA 程序具有
可扩展性，如图 \ref{fig:4-5} 所示。图中时间从上向下推进。在只有少量执行资源
的低成本系统中，同时执行的线程块较少，左侧示例每次执行两个线程块；在执行资源
更多的高端系统中，右侧示例每次执行四个线程块。如今的高端 GPU 可以同时执行
数百乃至数千个线程块。

同时执行的一组线程块通常称为一个\textbf{波次（wave）}。网格中的线程块总数
与能够同时执行的线程块数之比，称为波次数。例如，图 \ref{fig:4-5} 左侧示例
有四个波次，右侧示例有两个波次。当线程块之间负载不均衡时，较多波次更理想，
因为硬件有更多机会在 SM 之间平衡负载；当线程块负载均衡时，较少波次通常可以
接受，有时还更理想。第六章讨论线程粗化时会再次看到这一点。

\begin{figure}[htbp]
  \centering
  \scriptsize
  \begin{tabular}{c@{\hspace{1.0em}}c}
    \fbox{\begin{minipage}{0.34\textwidth}\centering
      \textbf{低成本 GPU：每次 2 个块}\\[0.2em]
      波次 1：块 0、1\\
      波次 2：块 2、3\\
      波次 3：块 4、5\\
      波次 4：块 6、7
    \end{minipage}}
    &
    \fbox{\begin{minipage}{0.34\textwidth}\centering
      \textbf{高端 GPU：每次 4 个块}\\[0.2em]
      波次 1：块 0、1、2、3\\
      波次 2：块 4、5、6、7
    \end{minipage}}
  \end{tabular}
  \caption{线程块之间没有同步约束时，CUDA 程序具有透明可扩展性
  （依据原书图 4.5 重绘）}
  \label{fig:4-5}
\end{figure}

使用较少波次时，需要避免最后一个波次不完整、只包含少量线程块而无法充分利用
硬件的情况。例如，一个有 660 个线程块的网格运行在能够同时执行 264 个线程块
的 GPU 上时，会产生
\[
  660/264=2.5
\]
个波次。如果线程块负载较均衡，先执行 264 个块，再执行 264 个块，最后执行
只有 132 个块的部分波次。这个部分波次会使 GPU 资源利用不足，这种现象有时称为
\textbf{尾部效应（tail effect）}。为避免尾部效应，最好把网格线程块数量配置为
同时可运行线程块数量的整数倍，从而得到整数个波次；该数量取决于 GPU 所拥有的
硬件资源。

能够在资源数量不同的各种 GPU 上执行同一份应用代码，使厂商可以按照不同市场
对成本、功耗和性能的要求生产多种 GPU。例如，移动处理器可以用极低功耗较慢地
执行应用，桌面处理器则可以用更高功耗更快地执行同一应用；两者不需要修改应用
代码。这种让同一应用代码在具有不同执行资源的硬件上运行的能力称为
\textbf{透明可扩展性（transparent scalability）}，它减轻了应用开发者的负担，
也提高了应用的可用性。

\section{Warp 与 SIMD 硬件}
\label{sec:warps-and-simd}

前面看到，当线程块处于不同的屏障同步范围时，它们可以按照任意相对顺序执行，
从而允许 CUDA 应用在不同设备之间透明扩展。但是，还没有讨论同一线程块内部线程
的执行时序。概念上，应假设线程块中的线程也可以按照任意相对顺序执行。在具有
多个阶段的算法中，只要希望屏障范围内的所有线程完成上一阶段后再进入下一阶段，
就应使用屏障同步。内核正确性不应依赖线程在没有屏障同步时恰好同步执行的假设。

CUDA GPU 中的线程调度是硬件实现概念，必须结合具体硬件讨论。在迄今为止的大多数
实现中，线程块被分配到流式多处理器后，会进一步划分为每组 32 个线程的单元，
称为 \textbf{warp}。warp 的大小由实现决定，未来 GPU 代际可能发生变化。了解
warp 有助于理解和优化特定 CUDA 设备代际上的应用性能。

warp 是 SM 中的线程调度单位。图 \ref{fig:4-6} 展示了一个实现中线程块划分为
warp 的方式。示例中，块 1、块 2 和块 3 都分配给同一个 SM；每个块又被划分为
多个 warp。每个 warp 由连续的 \code{threadIdx} 值组成：线程 0–31 构成第一个
warp，线程 32–63 构成第二个 warp，以此类推。如果每个线程块有 256 个线程，
则每个块有
\[
  256/32=8
\]
个 warp；SM 中有三个这样的块，因此共有 $8\times3=24$ 个 warp。

\begin{figure}[htbp]
  \centering
  \scriptsize
  \begin{tabular}{c|c|c}
    \toprule
    线程块 & 线程范围 & warp 数量\\
    \midrule
    Block 1 & 0--31,\ 32--63,\ $\ldots$,\ 224--255 & 8\\
    Block 2 & 0--31,\ 32--63,\ $\ldots$,\ 224--255 & 8\\
    Block 3 & 0--31,\ 32--63,\ $\ldots$,\ 224--255 & 8\\
    \midrule
    \multicolumn{2}{r|}{SM 中总计} & 24\\
    \bottomrule
  \end{tabular}
  \caption{线程块被划分为 warp 以进行线程调度（依据原书图 4.6 重绘）}
  \label{fig:4-6}
\end{figure}

如果线程块是一维数组，即只使用 \code{threadIdx.x}，划分很直接。一个 warp 中的
\code{threadIdx.x} 连续递增。warp 大小为 32 时，warp 0 从线程 0 到线程 31，
warp 1 从线程 32 到线程 63。一般而言，warp $n$ 从线程 $32n$ 开始，到线程
$32(n+1)-1$ 结束。如果线程块大小不是 32 的倍数，最后一个 warp 会用非活动线程
填充到 32 个位置。例如，一个有 48 个线程的块会划分为两个 warp，第二个 warp
用 16 个非活动线程填充。

对于包含多个线程维度的线程块，在划分 warp 之前，线程维度会投影为线性化的
行主序布局。线性布局中，较大的 y、z 坐标排在较小坐标之后。二维线程块的情况
是：先放置 \code{threadIdx.y=0} 的所有线程，再放置
\code{threadIdx.y=1} 的所有线程，以此类推；相同 y 值的线程按照递增的
\code{threadIdx.x} 连续排列。

\begin{figure}[htbp]
  \centering
  \scriptsize
  \begin{tabular}{c}
    \begin{tabular}{c|c|c|c}
      $T_{0,0}$ & $T_{0,1}$ & $T_{0,2}$ & $T_{0,3}$\\ \hline
      $T_{1,0}$ & $T_{1,1}$ & $T_{1,2}$ & $T_{1,3}$\\ \hline
      $T_{2,0}$ & $T_{2,1}$ & $T_{2,2}$ & $T_{2,3}$\\ \hline
      $T_{3,0}$ & $T_{3,1}$ & $T_{3,2}$ & $T_{3,3}$
    \end{tabular}\\[0.4em]
    $\longrightarrow\quad
    T_{0,0},T_{0,1},T_{0,2},T_{0,3},
    T_{1,0},T_{1,1},\ldots,T_{3,3}$
    \quad（行主序线性布局）
  \end{tabular}
  \caption{将二维线程放入线性布局（依据原书图 4.7 重绘）}
  \label{fig:4-7}
\end{figure}

图 \ref{fig:4-7} 上半部分是二维线程块，下半部分是线性化视图。前四个线程的
y 坐标为 0，按 x 坐标递增排列；接下来四个线程的 y 坐标为 1，也按 x 坐标递增
排列。这里的坐标分别对应 CUDA 字段 \code{threadIdx.y} 和 \code{threadIdx.x}。
示例中的 16 个线程构成半个 warp，另用 16 个非活动线程补足一个 32 线程 warp。

如果二维线程块为 $8\times8$，64 个线程会构成两个 warp。第一个 warp 从
$T_{0,0}$ 到 $T_{3,7}$，第二个 warp 从 $T_{4,0}$ 到 $T_{7,7}$。读者可以
自行画出该图验证这一点。

对于三维线程块，首先把 \code{threadIdx.z=0} 的所有线程放入线性顺序，其中
线程按二维线程块的方式排列；然后放入 \code{threadIdx.z=1} 的线程，以此类推。
例如，一个三维 $2\times8\times4$ 线程块（x 维为 4、y 维为 8、z 维为 2）
共有 64 个线程，会划分为两个 warp：第一个 warp 从 $T_{0,0,0}$ 到
$T_{0,7,3}$，第二个 warp 从 $T_{1,0,0}$ 到 $T_{1,7,3}$。

SM 设计为按照单指令多数据（Single Instruction, Multiple Data，SIMD）模型执行
warp 中的所有线程。也就是说，在任意时刻，一条指令会被取出并由 warp 中所有
线程执行。为此，一个 SM 被组织为多个处理分区，SM 上的 warp 分布到这些处理
分区中。图 \ref{fig:4-8} 给出了一个玩具示例：SM 由两个处理分区组成，每个
处理分区包含 8 个共享一条指令取出/分发单元的流式处理器。作为真实例子，拥有
128 个流式处理器的 Hopper H100 SM 被组织为四个处理分区，每个分区有 32 个
流式处理器。

\begin{figure}[htbp]
  \centering
  \scriptsize
  \begin{tabular}{c@{\hspace{1em}}c}
    \fbox{\begin{minipage}{0.31\textwidth}\centering
      \textbf{处理分区 0}\\
      指令取出/分发\\
      SP SP SP SP SP SP SP SP
    \end{minipage}}
    &
    \fbox{\begin{minipage}{0.31\textwidth}\centering
      \textbf{处理分区 1}\\
      指令取出/分发\\
      SP SP SP SP SP SP SP SP
    \end{minipage}}
  \end{tabular}\\[-0.1em]
  \scriptsize 一个玩具 SM；真实 H100 SM 有四个处理分区，每个分区 32 个 SP。
  \caption{SM 由处理分区组成，以执行 SIMD 指令（依据原书图 4.8 重绘）}
  \label{fig:4-8}
\end{figure}

同一 warp 中的线程会被分配到同一个处理分区，该分区为 warp 取出指令，并让
warp 中所有线程执行这条指令。由于 CUDA 线程调度机制把 warp 映射到处理分区，
实际上限制了 warp 中所有线程在任意时刻执行相同指令，因此 warp 的执行行为
通常称为单指令多线程（Single Instruction, Multiple Thread，SIMT）。

SIMT 的一个重要优点是，程序员可以用标量线程编写内核代码，之后由 SM 把线程
分组为 warp，透明地使用 SIMD 硬件。这不同于传统 CPU SIMD 编程：在 CPU 中，
程序员或编译器通常需要使用 API 在每个线程内部显式利用 SIMD 硬件，增加了另一
层编程复杂度。\footnote{虽然 SIMT 比 SIMD 更易编程，但程序员仍应避免控制流
分歧，因为分歧会对性能产生负面影响，见第 \ref{sec:control-divergence} 节。}

SIMD 的优势在于，指令取出/分发单元等控制硬件的成本由多个执行单元共享。这样
设计可以让较小比例的硬件用于控制，把更大比例的芯片面积用于提高算术吞吐量。
在可预见的未来，warp 划分仍会是常见的实现技术。不过，warp 大小可能因实现
而异；截至目前，所有 CUDA 设备都使用了类似的 32 线程 warp 配置。

由于同一 warp 中线程的调度方式具有特殊关系，CUDA 提供了一组 warp 级原语，即
在 warp 内进行高效数据交换和同步的 API 函数。程序员可以使用它们提高涉及线程
协作的计算速度和效率。这些原语支持以 warp 为中心的编程风格，把 warp 作为
软件层次结构中（线程块与线程之间）的额外层次。第十、十一和十二章将介绍相关
warp 级原语。

\begin{quote}
\textbf{Warp 与 SIMD 硬件}

1945 年，John von Neumann 在报告中描述了构建电子计算机的模型，该模型以先驱
EDVAC 计算机的设计为基础。今天通常称为“冯·诺依曼模型”，它几乎是所有现代
计算机的基础蓝图。

冯·诺依曼模型中的计算机有输入/输出（I/O），可以向系统提供程序和数据，也可以
从系统产生结果。执行程序时，计算机首先把程序及其数据输入内存。程序由一组
指令组成；控制单元维护程序计数器（Program Counter，PC），其中保存下一条待
执行指令的内存地址。在每个“指令周期”中，控制单元使用 PC 将指令取入指令
寄存器（Instruction Register，IR），再检查指令位以确定计算机各组件要执行的
动作。这也是该模型被称为“存储程序”模型的原因：用户只需把不同程序存入内存，
就能改变计算机行为。

将线程作为 warp 执行的动机，可以用适合 GPU 的修改版冯·诺依曼模型表示。对应
图 \ref{fig:4-8} 中处理分区的处理器只有一个负责取出和分发指令的控制单元；
相同的控制信号会发送到多个处理单元，每个处理单元对应 SM 中的一个流式处理器，
并执行 warp 中的一个线程。所有处理单元都由指令寄存器中的同一条指令控制，
它们的执行差异来自寄存器文件中不同的数据操作数。这就是处理器设计中的 SIMD。

例如，所有处理单元都可以执行 \code{add r1, r2, r3}，但不同处理单元中的
\code{r2} 和 \code{r3} 内容不同。SIMD 是 Flynn 分类法 \cite{flynn-4} 中四类
计算机架构之一，另外三类是单指令单数据（SISD）、多指令单数据（MISD）和多
指令多数据（MIMD）。现代处理器的控制单元相当复杂，包括指令取出逻辑和指令
缓存访问端口；让多个处理单元共享一个控制单元，可以显著降低硬件制造成本和
功耗。
\end{quote}

\section{控制流分歧}
\label{sec:control-divergence}

当 warp 中所有线程在处理数据时遵循相同的执行路径，更正式地说，遵循相同的
控制流时，SIMD 执行效果很好。例如，对于 \code{if-else} 结构，如果 warp 中
所有线程都执行 \code{then} 路径，或者所有线程都执行 \code{else} 路径，执行
就很高效。

但是，当同一 warp 中的线程采取不同控制流路径时，SIMD 硬件会分别遍历这些路径，
每条路径执行一次。例如，部分线程走 \code{then} 路径、其余线程走 \code{else}
路径时，硬件会执行两次：一次执行 \code{then} 路径，另一次执行 \code{else}
路径。每次遍历其中一条路径时，走另一条路径的线程都不会产生有效作用。

同一 warp 中的线程采取不同控制流路径时，称这些线程发生了\textbf{控制流分歧
（control divergence）}。这种多遍执行方式使 SIMD 硬件能够实现 CUDA 线程的
完整语义：硬件仍对 warp 中所有线程执行相同指令，但只让选择了当前路径的线程
在该遍中生效，从而既保持线程独立性，又利用 SIMD 硬件的低成本。分歧的代价是
额外路径遍历，以及每次遍历中非活动线程仍占用的执行资源。

\begin{figure}[htbp]
  \centering
  \scriptsize
  \begin{tabular}{c|c|c}
    \toprule
    线程 & \code{then} 遍 & \code{else} 遍\\
    \midrule
    0--23 & 执行 A & 非活动\\
    24--31 & 非活动 & 执行 B\\
    \bottomrule
  \end{tabular}\\[-0.1em]
  \scriptsize 两条路径执行完后，warp 中线程可能重新汇合并执行 C。
  \caption{warp 在 \code{if-else} 语句处分歧的示例（依据原书图 4.9 重绘）}
  \label{fig:4-9}
\end{figure}

图 \ref{fig:4-9} 中，包含线程 0–31 的 warp 到达 \code{if-else} 语句时，线程
0–23 走 \code{then} 路径，线程 24–31 走 \code{else} 路径。warp 首先执行一遍，
其中线程 0–23 执行 A、线程 24–31 非活动；随后再执行一遍，其中线程 24–31
执行 B、线程 0–23 非活动。之后线程可能重新汇合并执行 C。

在 Pascal 及更早架构中，这些遍按顺序执行，即一遍执行完后再执行另一遍。从
Volta 架构开始，不同遍可以并发执行，也就是一条路径的执行可以与另一条路径
交错。这一特性称为独立线程调度（independent thread scheduling）。因此，不能
假定 warp 中线程在分歧路径执行后一定会重新汇合。如果 warp 中所有线程必须
完成某个阶段后才能继续，应使用 \code{__syncwarp()} 等 warp 级屏障同步原语
保证正确性。

分歧也可能出现在其他控制流结构中。图 \ref{fig:4-10} 展示了分歧
\code{for} 循环：每个线程执行不同数量的迭代，次数从 4 到 8 不等。前四次迭代
中所有线程都活动并执行 A；第五次迭代中，部分线程执行 A，另一些线程（例如
线程 4）已经完成迭代而处于非活动状态；第八次迭代中，图示线程中只有线程 1
和线程 7 执行 A，其他线程均非活动。

\begin{figure}[htbp]
  \centering
  \scriptsize
  \begin{tabular}{c|cccccccc}
    \toprule
    迭代 & 1 & 2 & 3 & 4 & 5 & 6 & 7 & 8\\
    \midrule
    线程 1 & A&A&A&A&A&A&A&A\\
    线程 4 & A&A&A&A&--&--&--&--\\
    线程 7 & A&A&A&A&A&A&A&A\\
    \bottomrule
  \end{tabular}\\[-0.1em]
  \scriptsize A 表示执行，-- 表示非活动；不同线程的循环次数不同。
  \caption{warp 在 \code{for} 循环处分歧的示例（依据原书图 4.10 重绘）}
  \label{fig:4-10}
\end{figure}

检查控制结构的判定条件，可以判断它是否可能造成线程分歧。如果判定条件基于
\code{threadIdx} 值，控制语句就可能导致分歧。例如，条件
\code{threadIdx.x > 2} 会让块中第一个 warp 的线程走两条不同路径：
线程 0、1、2 走一条路径，线程 3、4、5 等走另一条路径。类似地，如果
循环条件依赖线程索引，循环也可能造成分歧。

使用带线程控制分歧的控制结构，一个常见原因是处理线程映射到数据时的边界条件。
这是因为线程总数通常必须是线程块大小的倍数，而数据大小可以是任意数值。从
第二章的向量加法内核开始，我们就使用了 \code{if (i < n)}，因为并非所有向量
长度都是线程块大小的倍数。

例如，假设向量长度为 1003，线程块大小为 64。需要启动 16 个线程块处理全部
1003 个元素，但 16 个线程块共有 1024 个线程，因此必须禁止第 15 个线程块中
最后 21 个线程执行原程序不允许的工作。这 16 个线程块被划分为 32 个 warp，
只有最后一个线程块中的第二个 warp 会产生控制流分歧。

随着处理向量长度增大，边界条件造成的分歧对性能的影响会降低。长度为 100 时，
四个 warp 中有一个分歧，影响可能很明显；长度为 1000 时，32 个 warp 中只有
一个分歧，分歧只影响约 3\% 的执行时间；即使使该 warp 执行时间加倍，对总时间
的影响也约为 3\%。长度达到 10,000 或更大时，只有 313 个 warp 中的一个分歧，
影响会远低于 1\%。

二维数据（例如第三章的彩色转灰度示例）也需要用 \code{if} 处理数据边缘线程
的边界条件。对 $62\times76$ 图片，使用 $4\times5=20$ 个 $16\times16$
线程块，每个块划分成 8 个 warp，每个 warp 包含两行线程，总计 160 个 warp。
参照第三章图 3.5：

\begin{itemize}
  \item 区域 1 的 12 个块中没有 warp 分歧，共 $12\times8=96$ 个 warp；
  \item 区域 2 的 24 个 warp 全部产生分歧；
  \item 区域 3 中，底部 warp 完全映射到图像外，因此不会通过 \code{if}，不产生分歧；
  \item 区域 4 中，前 7 个 warp 产生分歧，最后一个 warp 不产生分歧。
\end{itemize}

总计 160 个 warp 中有 31 个产生控制流分歧。如果处理 $200\times150$ 的图片，
使用 $16\times16$ 线程块，则共有 $13\times10=130$ 个线程块，即 1040 个 warp。
区域 1–4 的 warp 数分别为
\[
  864=12\times9\times8,\quad
  72=9\times8,\quad
  96=12\times8,\quad
  8=1\times8.
\]
其中只有 80 个 warp 产生分歧，性能影响小于 8\%。如果处理水平尺寸超过 1000
像素的真实图片，影响会低于 2\%。

\section{Warp 调度与延迟容忍}
\label{sec:warp-scheduling-latency-tolerance}

线程被分配到 SM 时，通常会有比 SM 中流式处理器更多的线程分配给它。也就是说，
每个 SM 的执行单元只能在任意时刻执行已分配线程的一部分。在较早的 GPU 设计中，
一个 SM 在任意时刻只能为一个 warp 执行一条指令；较新的设计中，一个 SM 可以
同时为少量 warp 执行指令，因为 SM 包含多个处理分区。无论哪种情况，硬件在
任意时刻都只能为分配给 SM 的 warp 子集执行指令。

为什么要给 SM 分配这么多 warp？答案是为了容忍全局内存访问等长延迟操作。当
warp 要执行的指令需要等待先前发起的长延迟操作结果时，该 warp 不会被选中执行；
硬件会选择另一个不再等待结果的驻留 warp。没有等待前序指令结果的 warp 已经
\textbf{准备好执行}。若有多个 warp 准备好执行，优先级逻辑会选择其中一个。

SM 在每个时刻选择不同准备好执行的 warp 的能力称为\textbf{细粒度多线程
（fine-grained multithreading）}。它让 SM 可以利用某些线程的指令等待时间，
执行其他线程的指令。因此，即使指令延迟很长，warp 的细粒度多线程执行仍能保持
执行单元高利用率。这通常称为\textbf{延迟容忍（latency tolerance）}或
\textbf{延迟隐藏（latency hiding）}。

\begin{quote}
\textbf{延迟容忍}

日常生活中经常需要容忍延迟。例如，在邮局寄包裹的人理想情况下应在到达服务柜台
前填好所有表格和标签。但现实中，有些人会等到柜员告诉他们填哪张表、如何填写。
柜台前排着长队时，应最大化柜员的工作效率。让一个人在柜台前填写表格、让所有
其他人等待，并不是好办法；柜员应在此人填写表格时服务下一位顾客。其他顾客
已经“准备好”，不应被需要更多时间填写表格的人阻塞。

因此，好的柜员会礼貌地请第一位顾客先到一旁填写表格，同时服务其他顾客。通常，
第一位顾客填完表格、柜员也服务完当前顾客后，就会立即得到服务，而不必重新排到
队尾。可以把这些顾客看作 warp，把柜员看作硬件执行单元；需要填写表格的顾客
对应于其后续执行依赖长延迟操作的 warp。
\end{quote}

warp 调度也可以容忍流水线浮点运算、条件分支等其他类型的操作延迟。只要有足够
多的 warp，硬件就很可能在任意时刻找到一个可执行 warp，使执行硬件保持充分利用，
同时让其他 warp 等待长延迟操作的结果。选择准备好的 warp 不会给执行时间线引入
空闲或浪费时间，这称为\textbf{零开销线程调度（zero-overhead thread scheduling）}。

\begin{quote}
\textbf{线程、上下文切换与零开销调度}

根据冯·诺依曼模型，现代计算机中的线程是一个程序及其在冯·诺依曼处理器上的
执行状态。线程包括程序代码、当前执行的指令，以及变量和数据结构的值。程序
代码保存在内存中，PC 跟踪正在执行的程序指令地址，IR 保存当前指令，寄存器
和内存保存变量与数据结构的值。

现代处理器支持上下文切换，让多个线程轮流推进。通过保存和恢复 PC、寄存器及
内存中的状态，可以暂停一个线程并在稍后正确恢复。然而，保存和恢复寄存器内容
会产生显著的执行时间开销。

零开销调度指 GPU 能够让等待长延迟指令结果的 warp 休眠，同时激活已经准备好的
warp，而不在处理单元中引入额外空闲周期。传统 CPU 的上下文切换会产生这种空闲
周期，因为切换线程需要把离开线程的寄存器内容保存到内存，再从内存加载进入线程
的执行状态。GPU SM 通过把已分配 warp 的所有执行状态保存在硬件寄存器中实现
零开销调度，因此在 warp 之间切换时不需要保存和恢复状态。
\end{quote}

warp 调度会把 warp 指令的长等待时间“隐藏”起来：硬件在等待期间执行其他 warp
的指令。容忍长操作延迟是 GPU 不像 CPU 那样把大量芯片面积用于缓存和分支预测
机制的主要原因。GPU 因而可以把更多芯片面积用于浮点执行单元和内存访问通道。

要使延迟容忍有效，最好给 SM 分配远多于其执行单元同时支持数量的线程。例如，
Hopper H100 GPU 的每个 SM 有 128 个流式处理器，但同时最多可以分配 2048 个
线程。因此，一个 SM 可以拥有多达
\[
  2048/128=16
\]
倍于当前一个时钟周期可执行数量的线程。这种线程相对于 SM 的过量配置
（oversubscription）对于延迟容忍非常重要，因为当前 warp 遇到长延迟操作时，
它提高了找到另一个 warp 执行的概率。

\section{资源划分与占用率}
\label{sec:resource-partitioning}

为了容忍较高的操作延迟，给 SM 分配很多 warp 是有利的；但 SM 不一定总能分配
它所支持的最大 warp 数。分配给 SM 的 warp 数与 SM 支持的最大 warp 数之比，
称为\textbf{占用率（occupancy）}。

除了流式处理器，SM 的执行资源还包括寄存器、共享内存（第五章讨论）、线程块
槽位和线程槽位。这些资源会在线程之间动态划分。例如，Hopper H100 GPU 每个
SM 最多支持 32 个线程块、64 个 warp（即 2048 个线程），每个线程块最多 1024
个线程。如果线程块大小为 1024，SM 的 2048 个线程槽位会分给 2 个线程块；
如果线程块大小分别为 512、256、128 或 64，则可分给 4、8、16 或 32 个线程块。

\begin{center}
\small
\begin{tabular}{r|r|r}
  \toprule
  线程块大小 & 每个 SM 可同时容纳的块数 & 线程总数\\
  \midrule
  1024 & 2 & 2048\\
  512 & 4 & 2048\\
  256 & 8 & 2048\\
  128 & 16 & 2048\\
  64 & 32 & 2048\\
  \bottomrule
\end{tabular}
\end{center}

线程槽位可以在不同线程块之间动态划分，使 SM 既能执行许多小线程块，也能执行
少量大线程块。这与固定划分形成对比：固定划分会在线程块需求小于固定资源时
浪费槽位，也无法支持需要更多槽位的线程块。

动态资源划分可能使不同资源限制之间产生微妙交互，造成资源利用不足。以 H100
为例，线程块大小从 1024 变化到 64 时，每个 SM 的线程块数从 2 变化到 32，
这些情况下分配给 SM 的线程总数都是 2048，占用率最大。但如果每个线程块只有
32 个线程，填满 2048 个线程槽位需要 64 个块；每个 H100 SM 只有 32 个块槽位，
最多只能容纳 32 个块，即
\[
  32\times32=1024
\]
个线程。此时占用率为
\[
  1024/2048=50\%.
\]
因此，要充分利用线程槽位并达到最大占用率，每个线程块至少需要 64 个线程。

另一种降低占用率的情况是线程槽位数不能被线程块大小整除。H100 每个 SM 最多
支持 2048 个线程；如果线程块大小为 768，SM 只能容纳 2 个线程块，即 1536 个
线程，剩余 512 个线程槽位未使用。此时占用率为
\[
  1536/2048=75\%.
\]

前面的讨论没有考虑寄存器和共享内存等资源约束。CUDA 内核中的自动变量会放入
寄存器；有些内核使用很多自动变量，有些只使用少量变量。因此，有些内核每线程
需要很多寄存器，有些需要很少寄存器。SM 动态划分寄存器后，可以在每线程寄存器
需求较少时容纳更多线程块，在需求较多时容纳更少线程块。

H100 GPU 每个 SM 最多有 65,536 个寄存器。若要达到满占用率，2048 个线程需要
的寄存器数意味着每线程不能超过
\[
  65{,}536/2{,}048=32
\]
个寄存器。如果内核每线程使用 64 个寄存器，65,536 个寄存器最多支持 1024 个
线程，因此无论线程块大小如何设置，占用率最多为 50\%。编译器有时可以通过寄存器
溢出减少每线程寄存器需求并提高占用率，但线程从内存访问溢出寄存器值通常会增加
执行时间，甚至使整个网格执行时间变长。

假设一个内核每线程使用 31 个寄存器，并配置为每块 512 个线程。SM 同时运行
\[
  2048/512=4
\]
个块，共使用
\[
  2048\times31=63{,}488
\]
个寄存器，低于 65,536 的上限。若再声明两个自动变量，使每线程使用 33 个
寄存器，则 2048 个线程需要 67,584 个寄存器，超过上限。CUDA 运行时可能只给
每个 SM 分配 3 个块，使寄存器需求降为
\[
  3\times512\times33=50{,}688.
\]
这样，SM 上的线程数从 2048 降为 1536，占用率从 100\% 降到 75\%。这种资源
使用略微增加却造成并行度和性能显著下降的现象，有时称为\textbf{性能悬崖
（performance cliff）} \cite{ryoo-4}。

所有动态划分资源之间会相互作用，因此准确确定每个 SM 上运行的线程数可能很难。
Occupancy Calculator \cite{occupancy-calculator-4} 可以根据特定设备和内核的资源
使用情况计算每个 SM 上实际运行的线程数；它是 NVIDIA Nsight Compute profiler
\cite{nsight-compute-4} 的一部分。主机代码也可以调用 CUDA Occupancy API。一个
常用的函数是
\[
\code{cudaOccupancyMaxActiveBlocksPerMultiprocessor()}
\]
它可以在给定 GPU 和网格配置时确定内核占用率。

\section{查询设备属性}
\label{sec:querying-device-properties}

前面对 SM 资源划分的讨论引出一个重要问题：如何知道某个设备可用的资源数量？
CUDA 应用运行时，如何知道设备中的 SM 数量，以及每个 SM 可以分配的线程块和
线程数量？应用通常需要在多种硬件系统上运行，因此需要查询底层硬件可用的资源
和能力，以便利用更强的系统，同时适配资源较少的系统。

\begin{quote}
\textbf{资源与能力查询}

日常生活中，我们也会查询环境的资源和能力。例如预订酒店时，可以查看房间配套
设施。若房间配有吹风机，就不必自己携带；若没有游泳池但有健身房，就可以准备
跑鞋和运动服。不同地区的酒店提供的牙膏、牙刷、洗发水、护发素、微波炉、冰箱
和泳池也可能不同。酒店设施就是酒店的属性，或者说资源与能力。经验丰富的旅行者
会在酒店网站上查询这些属性，选择更符合需求的酒店，并更有效地准备行李。
\end{quote}

每个 CUDA 设备 SM 的资源数量由设备的计算能力（compute capability）指定。一般
而言，计算能力级别越高，每个 SM 可用的资源越多，而且 GPU 的计算能力通常会
随代际增加。A100 GPU 的计算能力为 8.0，H100 GPU 的计算能力为 9.0。

CUDA C++ 提供了让主机代码查询系统中设备属性的内置机制。CUDA 运行时系统
（设备驱动）提供 \code{cudaGetDeviceCount} API，返回系统中 CUDA-capable
设备的数量。主机代码可以使用：

\begin{lstlisting}[language=C++,firstnumber=1]
int devCount;
cudaGetDeviceCount(&devCount);
\end{lstlisting}

现代 PC 往往有两个或更多 CUDA-capable 设备，因为许多 PC 带有一个或多个集成
GPU。集成 GPU 是默认图形单元，提供现代窗口用户界面所需的基本图形能力和硬件
资源，但大多数 CUDA 应用在这些设备上的性能并不理想。因此，主机代码可以遍历
所有可用设备，查询其资源和能力，并选择拥有足够资源、能够提供满意性能的设备。

CUDA 运行时把系统中的设备编号为 0 到 \code{devCount-1}。函数
\[
\code{cudaGetDeviceProperties}
\]
根据设备编号返回设备属性。例如，主机代码可以遍历设备：

\begin{lstlisting}[language=C++,firstnumber=1]
cudaDeviceProp devProp;
for (unsigned int i = 0; i < devCount; i++) {
    cudaGetDeviceProperties(&devProp, i);
    // Decide if device has sufficient resources/capabilities
}
\end{lstlisting}

\code{cudaDeviceProp} 是一个 C 结构体类型，其字段表示 CUDA 设备属性。CUDA C++
Programming Guide 列出了该类型的全部字段。调用完成后，属性会写入变量
\code{devProp}。

\code{devProp.maxThreadsPerBlock} 表示查询设备允许的每个线程块最大线程数。有些
设备允许每块最多 1024 个线程，有些设备允许更少；未来设备也可能允许更多。因此，
程序最好查询可用设备，判断它们是否允许应用所需的线程块大小。使用线程块集群的
CUDA 程序还应查询最大集群大小。可以调用
\[
\code{cudaOccupancyMaxPotentialClusterSize()}
\]
获取。

设备中的 SM 总数由 \code{devProp.multiProcessorCount} 给出。如果应用需要大量
SM 才能获得满意性能，应检查该属性。设备时钟频率保存在
\code{devProp.clockRate} 中。时钟频率与 SM 数量的组合，可以较好地表示设备的
最大硬件执行吞吐量。

主机代码可以通过以下字段查询线程块各维度允许的最大线程数：
\[
\begin{aligned}
&\code{devProp.maxThreadsDim[0]}\quad &&x\text{ 维},\\
&\code{devProp.maxThreadsDim[1]}\quad &&y\text{ 维},\\
&\code{devProp.maxThreadsDim[2]}\quad &&z\text{ 维}.
\end{aligned}
\]
例如，自动调优系统可以使用这些信息设置线程块维度的搜索范围。类似地，以下三个
字段
\[
\begin{aligned}
&\code{devProp.maxGridSize[0]},\quad
\code{devProp.maxGridSize[1]},\\
&\code{devProp.maxGridSize[2]}
\end{aligned}
\]
分别给出网格 x、y、z 维允许的最大线程块数。它们
可以用于判断一个网格是否有足够线程覆盖整个数据集，或者是否需要采用迭代方法。

\code{devProp.regsPerBlock} 给出单个线程块可使用的寄存器上限，可用于判断特定
设备上的内核是否能达到最大占用率，或者是否受寄存器使用限制。这个字段名称略有
误导：对于多数计算能力级别，一个线程块可使用的最大寄存器数确实等于 SM 中的
寄存器总数；但对于某些计算能力级别，一个线程块的寄存器上限小于 SM 的总量。

前文还讨论过 warp 大小取决于硬件。warp 大小可以通过
\code{devProp.warpSize} 获取。\code{cudaDeviceProp} 还有许多其他字段；随着
本书介绍相应概念和特征，我们会在后文讨论它们。

\section{小结}
\label{sec:chapter-4-summary}

CUDA-capable GPU 由流式多处理器（SM）组成，每个 SM 包含多个由流式处理器构成
的处理分区，并共享控制逻辑和内存资源。启动网格时，网格中的线程块以任意顺序
分配给 SM，从而使 CUDA 应用具有透明可扩展性。透明可扩展性带来一个约束：不同
线程块中的线程不应相互同步。从 Hopper 架构开始，SM 被分组为 GPC，线程块还可
选择分组为线程块集群。同一集群中的线程块在同一 GPC 上执行，可以相互同步并
访问彼此的分布式共享内存。

线程以线程块为单位分配给 SM。线程块分配到 SM 后，进一步划分为 warp。warp 中
的线程遵循单指令多数据（SIMD）模型执行。如果同一 warp 中的线程采取不同执行
路径，处理分区会分遍执行这些路径，每个线程只在与自己所选路径对应的遍中活动。

分配给 SM 的线程数可能远多于 SM 能同时执行的线程数。在任意时刻，SM 只执行
驻留 warp 子集中的指令，这让其他 warp 可以等待长延迟操作，而不影响大量执行
单元的总体吞吐量。分配给 SM 的线程数与 SM 支持的最大线程数之比称为占用率。
更高的占用率通常有助于隐藏长操作延迟并获得较高执行吞吐量。

不同 CUDA-capable 设备对每个 SM 可用资源的限制可能不同。例如，设备会限制每个
SM 可容纳的线程块数、线程数、寄存器数和其他资源数量。对于每个内核，一个或多个
资源限制可能成为占用率的决定因素。CUDA C++ 提供占用率计算器和 API，帮助开发者
为不同 GPU 选择合适的内核与启动配置；它也支持程序员在运行时查询 GPU 可用资源。

\phantomsection
\section*{练习}
\addcontentsline{toc}{section}{练习}

\begin{enumerate}[label=\arabic*.,leftmargin=2.7em]
  \item 考虑下面的 CUDA 内核及调用它的主机函数：

\begin{lstlisting}[language=C++,firstnumber=1]
__global__ void foo_kernel(int* a, int* b) {
    unsigned int i = blockIdx.x*blockDim.x + threadIdx.x;
    if (threadIdx.x < 40 || threadIdx.x >= 104) {
        b[i] = a[i] + 1;
    }
    if (i % 2 == 0) {
        a[i] = b[i] * 2;
    }
    for (unsigned int j = 0; j < 5 - (i % 3); ++j) {
        b[i] += j;
    }
}
void foo(int* a_d, int* b_d) {
    unsigned int N = 1024;
    foo_kernel<<<(N + 128 - 1)/128, 128>>>(a_d, b_d);
}
\end{lstlisting}

  \begin{enumerate}[label=\alph*.,leftmargin=2.2em]
    \item 每个线程块中有多少个 warp？
    \item 网格中有多少个 warp？
    \item 对第 04 行语句：
      \begin{enumerate}[label=\roman*.,leftmargin=2em]
        \item 网格中有多少个活动 warp？
        \item 网格中有多少个分歧 warp？
        \item 块 0 的 warp 0 的 SIMD 效率（百分比）是多少？
        \item 块 0 的 warp 1 的 SIMD 效率（百分比）是多少？
        \item 块 0 的 warp 3 的 SIMD 效率（百分比）是多少？
      \end{enumerate}
    \item 对第 07 行语句：
      \begin{enumerate}[label=\roman*.,leftmargin=2em]
        \item 网格中有多少个活动 warp？
        \item 网格中有多少个分歧 warp？
        \item 块 0 的 warp 0 的 SIMD 效率（百分比）是多少？
      \end{enumerate}
    \item 对第 09 行循环：
      \begin{enumerate}[label=\roman*.,leftmargin=2em]
        \item 有多少次迭代没有分歧？
        \item 有多少次迭代产生分歧？
      \end{enumerate}
  \end{enumerate}

  \item 对向量加法，假设向量长度为 2000，每个线程计算一个输出元素，线程块
  大小为 512。网格中有多少个线程？

  \item 对上一题，考虑长度为 2000 的边界检查，预计有多少个 warp 会产生分歧？

  \item 假设一个有 8 个线程的线程块在到达屏障前执行一段代码。8 个线程执行
  这段代码所需时间（微秒）分别为 2.0、2.3、3.0、2.8、2.4、1.9、2.6 和
  2.9，其余时间都在等待屏障。线程总执行时间中有百分之多少用于等待屏障？

  \item CUDA 程序员认为，如果每个线程块只启动 32 个线程，那么在需要屏障同步
  的地方可以省略 \code{__syncthreads()}。你认为这是好主意吗？请解释原因。

  \item 如果 CUDA 设备的 SM（流式多处理器）最多容纳 1536 个线程和 4 个线程块，
  下列哪一种线程块配置会使 SM 中的线程数最多？
  \begin{enumerate}[label=\alph*.,leftmargin=2.2em]
    \item 每块 128 个线程；
    \item 每块 256 个线程；
    \item 每块 512 个线程；
    \item 每块 1024 个线程。
  \end{enumerate}

  \item 假设设备每个 SM 最多允许 64 个线程块和 2048 个线程。指出下列每种
  每个 SM 的分配是否可行；若可行，给出占用率：
  \begin{enumerate}[label=\alph*.,leftmargin=2.2em]
    \item 8 个块，每块 128 个线程；
    \item 16 个块，每块 64 个线程；
    \item 32 个块，每块 32 个线程；
    \item 64 个块，每块 32 个线程；
    \item 32 个块，每块 64 个线程。
  \end{enumerate}

  \item 某 GPU 的硬件限制为：每个 SM 2048 个线程、32 个线程块和 64 K（65,536）
  个寄存器。对下列内核特征，判断内核能否达到满占用率；如果不能，指出限制因素：
  \begin{enumerate}[label=\alph*.,leftmargin=2.2em]
    \item 每块 128 个线程，每线程 30 个寄存器；
    \item 每块 32 个线程，每线程 29 个寄存器；
    \item 每块 256 个线程，每线程 34 个寄存器。
  \end{enumerate}

  \item 一名学生说，他用每块 $32\times32$ 个线程的矩阵乘法内核成功计算了两个
  $1024\times1024$ 矩阵的乘法。该学生使用的 CUDA 设备每块最多允许 512 个线程、
  每个 SM 最多允许 8 个线程块；并且每个线程计算结果矩阵的一个元素。你对此有
  什么反应？为什么？
\end{enumerate}

\begin{chapterbibliography}{9}
  \bibitem{cuda-guide-4}
  NVIDIA, \emph{CUDA C++ Programming Guide}, Release 12.3, 2023.
  \url{https://docs.nvidia.com/cuda/cuda-c-programming-guide/}.
  \bibitem{flynn-4}
  M. Flynn, ``Very high-speed computing systems,'' \emph{Proceedings of the IEEE},
  54(12), 1966, pp. 1901--1909.
  \bibitem{volta-4}
  NVIDIA, \emph{NVIDIA Tesla V100 GPU Architecture}, Version
  WP-08608-001\_v1.1, 2017.
  \bibitem{ryoo-4}
  S. Ryoo, C. I. Rodrigues, S. S. Stone, S. S. Baghsorkhi, S.-Z. Ueng,
  J. A. Stratton, W.-m. W. Hwu, ``Program optimization space pruning for a
  multithreaded GPU,'' in \emph{Proceedings of the 6th Annual IEEE/ACM
  International Symposium on Code Generation and Optimization}, 2008,
  pp. 195--204.
  \bibitem{occupancy-calculator-4}
  NVIDIA, \emph{Nsight Compute: Occupancy Calculator}, 2023.
  \url{https://docs.nvidia.com/nsight-compute/NsightCompute/index.html#occupancy-calculator}.
  \bibitem{nsight-compute-4}
  NVIDIA, \emph{Nsight Compute}, 2023.
  \url{https://docs.nvidia.com/nsight-compute/index.html}.
\end{chapterbibliography}
